\section{指标定理的陈述}
% ===== 原书 p.14 =====
\subsection{象征映射与椭圆复形}
设紧空间$X$上有椭圆算子$D:\Gamma E\to\Gamma F$。于是对每个$x\in X$与$\xi\in TX_x$，$D$的象征$\sigma(x,\xi)$是同态$E_x\to F_x$，它随$x$、$\xi$光滑变化。可见这一族同态给出一个齐次复形
\[
0\longrightarrow\pi^*E\xrightarrow{\ \alpha\ }\pi^*F\longrightarrow 0,
\]
其底空间为$TX$。\footnote{本节与下一节用$TX$表示$X$的余切丛（cotangent bundle）。注意若在$X$上取定度量，则可把$X$的切丛与余切丛等同。}由于$D$椭圆，已知$\xi\neq 0$时$\sigma(x,\xi)$可逆，故该复形在零截面之外正合。可见这个复形定义$K(TX)$中的一个元素，称为$D$的象征，记作$\sigma(D)$。

% ===== 原书 p.15 =====
这一构造几乎不加改动即可推广到微分算子复形。给定紧空间$X$上的一列向量丛$E^0,\ldots,E^n$与偏微分算子$d_i:\Gamma E^i\to\Gamma E^{i+1}$，称复形
\[
0\longrightarrow\Gamma E^0\xrightarrow{\ d_0\ }\Gamma E^1\xrightarrow{\ d_1\ }\cdots\xrightarrow{\ d_{n-1}\ }\Gamma E^n\longrightarrow 0
\]
为椭圆复形，若$d_{i+1}d_i=0$，且相应的象征复形
\[
0\longrightarrow\pi^*E^0\xrightarrow{\ \sigma_0\ }\pi^*E^1\xrightarrow{\ \sigma_1\ }\cdots\xrightarrow{\ \sigma_{n-1}\ }\pi^*E^n\longrightarrow 0
\]
在$TX$上于零截面之外正合。与前面一样，这个复形给出一个元素$\sigma(D^\bullet)\in K(TX)$，称为椭圆复形$D^\bullet$的象征。对椭圆复形，取代指标的是欧拉示性数（Euler characteristic），它由
\[
\chi(D^\bullet)=\sum_{i=0}^n(-1)^i\dim H^i(D^\bullet)
\]
定义，其中$H^i(D^\bullet)=\Ker d_i/\operatorname{Im} d_{i-1}$是复形的上同调。注意当复形$D^\bullet$只是单个算子时，欧拉示性数化为前面关于指标的定义。进一步，若按§1构造相应的长度为 1 的复形，便得到椭圆算子
\[
D:\Gamma\Big(\bigoplus_iE^{2i}\Big)\to\Gamma\Big(\bigoplus_iE^{2i+1}\Big),
\]
它由$D=\bigoplus_i d_{2i}+\bigoplus_i d^*_{2i+1}$给出，其中$d^*$表示关于某个度量的伴随。于是注意$\sigma(D)=\sigma(D^\bullet)\in K(TX)$且$\Ind D=\chi(D^\bullet)$，问题就化归为单个椭圆算子的情形。

\begin{example}
考虑（复化后的）de Rham 复形
\[
0\longrightarrow\Gamma\Lambda^0(TX\otimes\C)\xrightarrow{\ d\ }\Gamma\Lambda^1(TX\otimes\C)\xrightarrow{\ d\ }\cdots\xrightarrow{\ d\ }\Gamma\Lambda^n(TX\otimes\C)\longrightarrow 0,
\]
其中$\Gamma\Lambda^i(TX)$是$X$上次数为$i$的微分形式空间，$d$是外微分（exterior derivative）。此时复形的上同调就是 de Rham 上同调，$H^i(D^\bullet)=H^i_{\mathrm{dR}}(X)\otimes\C$，于是 de Rham 复形的欧拉示性数等于流形$X$的欧拉示性数：
\[
\chi(D^\bullet)=\sum_i\dim H^i_{\mathrm{dR}}(X)=\sum_i(-1)^i\dim H^i(X;\R)=\chi(X).
\]
相应的元素$\sigma(D^\bullet)\in K(TX)$称为 de Rham 象征，记作$\rho_X$。
\end{example}

\transnote{原文第一项写作$\sum_i\dim H^i_{\mathrm{dR}}(X)$，缺少$(-1)^i$。按欧拉示性数的定义应为$\sum_i(-1)^i\dim H^i_{\mathrm{dR}}(X)$；此处照译保留，读时须补符号。}

\subsection{拓扑指标的构造}
由于象征给出从$X$上椭圆算子到$K(TX)$的映射，我们希望把它与适当的映射$K(TX)\to\Z$复合。最先想到的一个，是点包含$i:\pt\to TX$诱导的映射。然而，回想定义$K(TX)=\tilde K(TX^+)$，会失望地看到映射$i^*:K(TX)\to K(\pt)\cong\Z$恒为零。事实上，注意$i$延拓为映射$i:\pt\to TX^+$，而约化$K$群$\tilde K(TX^+)$按定义是$\Ker i^*$，上述结论立即得到。所以必须到更深处寻找所需的映射。

% ===== 原书 p.16 =====
与其考虑$X$中基点的包含，不如整体地思考。回想若$X$是紧流形，则可把$X$嵌入适当大的$\R^n$\footnote{这可用涉及单位分解的简单构造说明。}。给定这样的嵌入$X\subset\R^n$，可把它延拓为切丛的嵌入$TX\subset T\R^n$。由 Thom 同构定理$K(T\R^n)\cong\Z$，故只要构造出映射$K(TX)\to K(T\R^n)$就满意了。

更一般地，若$i:X\to Y$是紧子流形$X$到$Y$的包含，可把它延拓为切丛的包含$i:TX\to TY$。随后如下构造映射
\[
i_!:K(TX)\to K(TY)
\]
（注意该映射关于包含是函子性的）。取定$Y$上的度量，令$N$是$X$在$Y$中的管状邻域，它与$X$在$Y$中的法丛微分同胚。延拓到切丛，可见$TN$也是$TX$在$TY$中的管状邻域，与$TX$在$TY$中的法丛微分同胚。管状邻域定理的证明见[16, §11]。由于$TN$是$TX$上的实向量丛，我们希望给它一个复结构，以便应用 Thom 同构。

考虑$X$上任一向量丛$E$，回想$E$局部同胚于$U\times\R^m$，其中$U\subset X$同胚于$\R^n$的一个开区域。于是$E$在$(x,\xi)\in E$处的切空间为$TE_{(x,\xi)}=TU_x\times T\R^m_\xi\cong TX_x\oplus E_x$，可见$E$上的切丛$TE$有分解
\[
TE=\pi^*TX\oplus\pi^*E,
\]
其中$\pi:E\to X$是投影。特别地，取$E=TX$与$E=TY|_X$，得到
\begin{align*}
T(TX)&\cong\pi^*TX\oplus\pi^*TX,\\
T(TY)|_{TX}&\cong\pi^*(TY|_X)\oplus\pi^*(TY|_X),
\end{align*}
其中$T(TX)$与$T(TY)|_{TX}$都是$TX$上的丛，$\pi:TX\to X$是投影。又因$N$与$TN$分别是$X$在$Y$中、$TX$在$TY$中的法丛，有$N\oplus TX\cong TY|_X$与$TN\oplus T(TX)\cong T(TY)|_{TX}$。于是$TN$分解为
\[
TN\cong\pi^*N\oplus\pi^*N.
\]
这样就能给$TN$赋予复向量丛的结构：
\[
TN\cong\pi^*N\oplus i\pi^*N\cong(\pi^*N)\otimes_\R\C\cong\pi^*(N\otimes_\R\C).
\]
于是把 Thom 同构$\varphi:K(TX)\to K(TN)$与自然的延拓同态$h:K(TN)\to K(TY)$复合，便得到所需映射
\[
i_!:K(TX)\xrightarrow{\ \varphi\ }K(TN)\xrightarrow{\ h\ }K(TY).
\]
注意该映射不依赖邻域$N$的选取，且是函子性的（即$(j\circ i)_!=j_!\circ i_!$），因为 Thom 同构传递。还有
\[
i^*\circ i_!:x\mapsto i^*\lambda_{TN}\cdot x=\Big(\sum_{i=0}^n(-1)^i\Lambda^i(N\otimes_\R\C)\Big)\cdot x.
\]

% ===== 原书 p.17 =====
\begin{definition}
若$X$是紧流形，取充分大的$m$与嵌入$i:X\to\R^m$，并令$j:P\to\R^m$是原点的包含。注意$j_!:K(TP)\to K(T\R^m)$就是 Thom 同构$\varphi:K(\pt)\to K(\C^m)$，故可把拓扑指标定义为复合
\[
\tind:K(TX)\xrightarrow{\ i_!\ }K(T\R^m)\xrightarrow{\ j_!^{-1}\ }K(TP)\cong\Z.
\]
\end{definition}

设已有两个嵌入$i:X\to\R^m$与$i':X\to\R^n$。考虑对角嵌入$i\oplus i':X\to\R^{m+n}$，可见$i\oplus i'$同痕（isotopic）于嵌入$i\oplus 0$与$0\oplus i'$。令$\lambda_n\in K(\C^n)$表示$\C^n\cong T\R^n$的 Thom 类，注意$(i\oplus 0)_!(x)=i_!(x)\cdot\lambda_n$与$(0\oplus i')_!(x)=\lambda_m\cdot i'_!(x)$。于是由 Thom 同构的传递性，$i\oplus 0$与$i$确定同一个 t-ind，$0\oplus i'$与$i'$亦然。可见拓扑指标不依赖嵌入的选取。

现已构造出所需映射$\tind:K(TX)\to\Z$，可把它与象征映射$\sigma:(\text{椭圆算子})\to K(TX)$复合。终于可以陈述本文的主要结果。

\begin{theorem}[Atiyah–Singer 指标定理]
若$D$是紧流形上的椭圆算子，则
\[
\Ind D=\tind\sigma(D).
\]
\end{theorem}

\transnote{$i^*:K(TX)\to K(\pt)$恒为零，是因为$\tilde K(TX^+)$就定义为$\Ker i^*$。所以指标不能靠「在一点处取纤维」读出；第 2.2 节绕开的正路是嵌入$\R^m$、把法丛复化，再用 Thom 同构造出$\tind$。}

\subsection{指标定理的上同调形式}
在进入 Atiyah–Singer 指标定理的$K$理论证明之前，先短暂离题，用示性类与上同调重新表述该定理。1963 年 Atiyah 与 Singer 最初证明时用的正是这种形式，工具来自配边理论。$K$理论表述或许能给出更优雅、更一般的证明，而上同调形式便于直接计算例子。为用上同调写出拓扑指标，设$X$是嵌入$\R^k$的紧$n$维流形，考虑如下非交换图：
\[
\begin{tikzcd}[column sep=3.4em,row sep=2.8em]
K(TX)\arrow[r,"\varphi"]\arrow[d] & K(TN)\arrow[r,"h"]\arrow[d] & K(T\R^k)\arrow[d] & K(TP)\cong\Z\arrow[l,"\varphi'"']\arrow[d]\\
H^*_c(TX)\arrow[r,"\psi"] & H^*_c(TN)\arrow[r,"k"] & H^*_c(T\R^k) & H^*_c(TP)\cong\Q\arrow[l,"\psi'"']
\end{tikzcd}
\]
其中$\varphi$、$\psi$、$\varphi'$、$\psi'$是各种 Thom 同构，$h$与$k$是延拓同态，竖直映射是 Chern 特征$\ch:K(\cdot)\to H^{\mathrm{even}}_c(\cdot;\Q)$。此图虽然不交换，但回想§1：对$X$上的复向量丛$E$，$K$理论与上同调的 Thom 同构由下式相联系：
\[
\psi^{-1}\circ\ch\circ\varphi:x\mapsto(-1)^n\Td(\overline{E})^{-1}\smile\ch(x).
\]

此处因$\overline{T\R^k}$显然是点$TP$上的平凡丛，有$\Td(\overline{T\R^k})=1$。又因$TN\cong\pi^*N\otimes_\R\C$，可见$TN\cong\overline{TN}$。注意$T(TX)\oplus TN=T(T\R^k)|_{TX}$是$TX$上的平凡丛，且$T(TX)\cong\pi^*TX\oplus\pi^*TX\cong\pi^*TX\otimes_\R\C$，得到
\[
\Td(\overline{TN})^{-1}=\Td(TN)^{-1}=\Td(T(TX))=\pi^*\Td(TX\otimes_\R\C)=\pi^*\IndexClass(X),
\]
% ===== 原书 p.18 =====
其中最后一个等式给出流形$X$的指标类$\IndexClass(X)$的定义。于是类$\Td(\overline{TN})^{-1}=\pi^*\IndexClass(X)$与嵌入无关。因此
\begin{align*}
\psi'^{-1}\circ\ch\circ\varphi'&:x\mapsto(-1)^k\Td(\overline{T\R^k})^{-1}\smile\ch(x)=(-1)^k\ch(x),\\
\psi^{-1}\circ\ch\circ\varphi&:x\mapsto(-1)^{n-k}\Td(\overline{TN})^{-1}\smile\ch(x)=(-1)^{n-k}\IndexClass(X)\smile\ch(x).
\end{align*}
回想$\tind=\varphi'^{-1}\circ h\circ\varphi$，经一番扩展的图追踪得到
\[
\ch\circ h\circ\varphi(x)=\ch\circ\varphi'(\tind x)=(-1)^k\psi'(\tind x)=(-1)^k(\tind x)\psi'(1),
\]
其中$\psi'(1)$是$H^*_c(T\R^k)$的典范生成元。再令$[T\R^k]$为$T\R^k$的基本同调类（fundamental homology class），有$\psi'(1)[T\R^n]=1$，于是算出
\begin{align*}
\tind x&=(\tind x\smile\psi'(1))[T\R^n]=(-1)^k(\ch\circ h\circ\varphi(x))[T\R^n]\\
&=(-1)^k(\ch\circ\varphi(x))[TN]=(1)^n\big(\psi(\ch(x)\smile\IndexClass(X))\big)[TN]\\
&=(-1)^n(\ch(x)\smile\IndexClass(X))[TX].
\end{align*}

\transnote{上式中的$[T\R^n]$照原书排印。本段的类都是$T\R^k$的基本类，故$[T\R^n]$应读作$[T\R^k]$：$n$是流形$X$的维数，$k$是$X$所嵌入的$\R^k$的维数。}

现在陈述 Atiyah–Singer 指标定理的拓扑版本。

\begin{theorem}[指标定理 A]
若$D$是紧$n$维流形$X$上的椭圆算子，则
\[
\Ind D=(-1)^n\big(\ch(\sigma(D))\smile\IndexClass(X)\big)[TX],
\]
其中$\IndexClass(X)=\Td(TX\otimes_\R\C)\in H^*(X;\Q)$是$X$的指标类。
\end{theorem}

为得到 1963 年指标定理论文中的确切陈述，注意若$X$是定向流形，则对每个$u\in H^n(X;\Z)$有$\psi(u)[TX]=(-1)^{n(n-1)/2}u[X]$，其中$\psi:H^*(X)\to H^*(TX)$是 Thom 同构。因我们对$X$的定向作了选取，必须引入符号$(-1)^{n(n-1)/2}$。于是得到原来的指标定理：

\begin{theorem}[指标定理 B]
若$D$是紧定向$n$维流形$X$上的椭圆算子，则
\[
\Ind D=(-1)^{n(n+1)/2}\big(\psi^{-1}\circ\ch(\sigma(D))\smile\IndexClass(X)\big)[X],
\]
其中$\IndexClass(X)=\Td(TX\otimes_\R\C)\in H^*(X;\Q)$是$X$的指标类。
\end{theorem}

为便于对具体的椭圆算子或复形计算拓扑指标，下面给出 Atiyah–Singer 指标定理的一个特例。在适当情形下，可化简定理陈述中的$\psi^{-1}\circ\ch(\sigma(D))$项。令$\psi:H^*(X)\to H^*(TX)$是 Thom 同构，$i^*:H^*(TX)\to H^*(X)$是零截面$X\to TX$诱导的映射，回想§1：$i^*\psi(x)=x\smile e(X)$，其中$e(X)=e(TX)\in H^n(X;\Z)$是$X$的 Euler 类。于是对一切$y\in H^*(TX)$有$i^*(y)=\psi^{-1}(y)\smile e(X)$，并得到
\[
\psi^{-1}\circ\ch(x)\smile e(X)=i^*(\ch(x))=\ch(i^*(x)).
\]
再设$D^\bullet$是紧流形$X$上的椭圆复形
\[
0\longrightarrow\Gamma E^0\xrightarrow{\ d_0\ }\Gamma E^1\xrightarrow{\ d_1\ }\cdots\xrightarrow{\ d_{n-1}\ }\Gamma E^n\longrightarrow 0,
\]
可见$i^*\sigma(D^\bullet)=\sum_{i=0}^n(-1)^iE^i$，于是
\[
\psi^{-1}\circ\ch(\sigma(D^\bullet))\smile e(X)=\ch\Big(\sum_{i=0}^n(-1)^iE^i\Big).
\]
现在把这个结果并入 Atiyah–Singer 指标定理的陈述。

\transnote{计算链中$(1)^n$一项照原书排印。由上一段的$\ch\circ\varphi(x)=(-1)^{n-k}\psi(\ch(x)\smile\IndexClass(X))$，跨到$[TN]$时应得$(-1)^k\cdot(-1)^{n-k}=(-1)^n$，故此处读作$(-1)^n$才与末行一致。}
